\section{Related Work}\label{sec:related-work}

\paragraph{Benchmarks from repository history.}
SWE-bench builds tasks from real GitHub issues and the pull requests that resolved them, and grades a patch by the tests that accompany the fix~\cite{jimenez2024swebench}. We borrow the idea of replaying real history: our feature tasks start at a commit's parent, give the agent a written ticket, and are graded by that commit's own tests. We differ in scope and purpose. Our tasks come from one private repository, not many public ones, and we vary the agent's model and effort rather than comparing systems on a leaderboard. We also add task kinds that history does not supply directly, namely locate and review tasks, and we validate each task with an oracle self-test.

\paragraph{Task-difficulty horizons.}
Kwa et al.\ measure the length of a task, in human working time, at which a model succeeds half the time~\cite{kwa2025longtasks}. We use the same shape of idea on a smaller scale. Our code-trace family is graded by program length in statements, not by human time, and we report the longest length each config solved rather than a fitted 50\,\% horizon, because one trial per cell does not support a fit. Our coding tasks have no clean size axis, so we cannot estimate a horizon for them at all.

\paragraph{Measurement models and small-sample intervals.}
The Rasch model expresses the chance of success as a function of the gap between a person's ability and an item's difficulty~\cite{rasch1960probabilistic}. We propose a version in which a config's ability combines a model term, an effort term, and a model-by-task term, and we state what the current data can and cannot estimate; most of it cannot yet be fitted because nearly every cell passed. Because our cells hold few trials and pass rates sit near 100\,\%, we report Wilson score intervals~\cite{wilson1927probable}, which behave well in that regime. The intervals are wide, and we say so wherever a comparison depends on them.

\paragraph{Planted faults.}
Mutation testing plants faults in a program to judge how well a test suite detects them~\cite{demillo1978hints}. We plant faults for a different reader: to give an agent a fix task with a known cause, or a review task with a known answer key. Our fix tasks also place the cause at different distances from the symptom, and some hide the test that would catch it. Unlike mutants generated by fixed operators, our planted bugs were written by a model, so their validity rests on an oracle check and spot reads, not on human review.

\paragraph{Model routing and reasoning effort in agent workflows.}
Orchestrating agents delegate to sub-agents, and users choose a model and an effort level for each. A common rule of thumb is a small model for lookups, a larger one for hard problems, and high effort when a task seems difficult. We know of no published measurement that tests such rules or compares effort against a second model at equal cost. The gap we address is concrete. We measure pass rate, dollars, and agent time for the same tasks across models and effort levels, run the agents in a setting close to real orchestration, and turn the results into a routing table whose rules we label by how well the evidence supports them.
